\section{METHODS}

\paragraph{Data Collection.}
Two mutually independent datasets were employed to train and validate the predictive models, thereby rigorously appraising architectural generalizability and achieving robust cytotoxicity prediction for selenium nanoparticles (SeNPs).

\begin{enumerate}
  \item \textbf{Generalized Benchmark Dataset}:
This dataset was derived from the meta-analysis of nanoparticle cytotoxicity conducted by Masarkar et al. \cite{masarkar2024}.
It compiles published experimental cytotoxicity records spanning four major nanomaterial categories: metals, metal oxides, polymers, and carbon-based nanoparticles, representing one of the benchmark public repositories in machine-learning-driven nanotoxicity prediction.
Because our architectural extensions build directly upon their established framework, this benchmark dataset was utilized to reproduce baseline models, thereby guaranteeing rigorous comparability and procedural validity during model optimization.
The raw benchmark dataset originally comprised 4,111 experimental entries, from which 4,102 curated records were retained following quality screening for subsequent train--test partitioning and algorithmic benchmarking.

  \item \textbf{Curated Domain-Specific Dataset}:
Specifically tailored to selenium nanoparticles, a specialized cytotoxicity dataset was systematically constructed by curating experimental records from 43 peer-reviewed publications \cite{chen_selenium_2008, li_reversal_2011, zheng_sialic_2011, liu_selenium_2012, wu_surface_2012, chen_surface_2012,
  yu_positive_2012, chen_peg-nanolized_2012, pi_pathway_2013, wu_induction_2013, feng_differential_2014,
  huang_selective_2013, li_functionalized_2013, zhang_enhancement_2013, jiang_gracilaria_2014, yu_ph-responsive_2014,
  zeng_construction_2015, liu_sequentially_2016, liu_facile_2016, yu_x-ray-responsive_2016, fu_rgd_2016,
  fang_targeting_2017, chen_efficient_2018, fang_dualtargeted_2018, huang_autophagy_2018, liu_highly_2018,
  wang_regulating_2016, hu_selenium_2019, zeng_potentiation_2019, huang_dual-targeting_2019, huang_highly_2019,
  zhang_systematic_2019, zheng_thermosensitive_2019, huang_chiralitydriven_2020,
  yang_lentinan-functionalized_2020, liu_rapid_2021, an_functionalized_2021, song_translational_2021,
  yang_atherosclerotic_2022, yang_functionalized_2021, liang_translational_2022, liu_specific_2022, du_highly_2022}.
We established a hybrid data curation protocol coupling coarse automated extraction via large language models (LLMs) with rigorous manual verification.
Specifically, prompt engineering was employed to direct an LLM in coarsely extracting feature attributes aligned with the schema defined by Masarkar et al. \cite{masarkar2024} across the 43 publications. For data points not explicitly reported in the main text but accessible in tables or figures, precise manual digitizing and extraction were conducted. Finally, all compiled records were cross-checked and validated against the primary literature by domain researchers.
The resulting curated dataset comprises 383 validated entries, each parameterized by 37 features spanning SeNP physicochemical descriptors, exposure conditions, and biological endpoints, thereby providing a comprehensive representation of structure--cytotoxicity relationships in selenium nanostructures.
This dataset provides the core empirical foundation for training, validating, and evaluating our enhanced predictive architectures.
\end{enumerate}

To elucidate the compositional landscape of the curated dataset, the 43 source publications were categorized into four thematic clusters based on their titles and primary investigative scopes: Material Synthesis/Surface Functionalization, Mechanistic Toxicology/Biosafety, Drug Delivery/Combination Therapy, and Theranostics/Immunotherapy/Other Applications. This taxonomic stratification serves solely to demonstrate literature coverage and thematic breadth, without altering the extraction criteria for experimental records in downstream modeling. Figure~\ref{fig:senp-literature-cloud} depicts the publication year point cloud and categorical distribution of the source literature.

\begin{figure}[H]
  \centering
  \begin{tikzpicture}
    \begin{axis}[
      acs axis,
      width=0.82\linewidth,
      height=6.5cm,
      xmin=2007.5,
      xmax=2022.8,
      ymin=0.55,
      ymax=4.45,
      xlabel={Publication year},
      xtick={2008,2010,2012,2014,2016,2018,2020,2022},
      ytick={1,2,3,4},
      yticklabels={Material/Surface,Mechanism/Safety,Delivery/Combination,Theranostics/Immuno/Other},
      y tick label style={font=\scriptsize},
      title={\textbf{(A)} Publication point cloud ($N = 43$)}
    ]
      \addplot+[only marks, mark=*, mark size=1.9pt, acsblue] coordinates {
        (2008,0.78) (2011,1.10) (2011.88,0.92) (2012.12,1.24)
        (2012.36,0.80) (2013.16,1.08) (2014,0.92) (2019,1.20)
      };
      \addplot+[only marks, mark=square*, mark size=1.8pt, acsred] coordinates {
        (2011.15,1.80) (2012.86,2.08) (2013.05,1.86)
        (2013.30,2.22) (2018,1.92) (2019.15,2.12)
      };
      \addplot+[only marks, mark=triangle*, mark size=2.1pt, acsorange] coordinates {
        (2012.44,2.76) (2012.68,3.18) (2013.42,2.90) (2014.04,3.24)
        (2014.32,2.78) (2015,3.12) (2015.96,2.86) (2016.22,3.22)
        (2017,2.94) (2018.16,3.16) (2019.30,2.84) (2020,3.12)
        (2021,2.90) (2021.24,3.20)
      };
      \addplot+[only marks, mark=diamond*, mark size=2pt, acsgray] coordinates {
        (2016.34,3.78) (2016.54,4.06) (2016.74,4.30) (2018.30,3.88)
        (2018.50,4.18) (2019.44,3.78) (2019.64,4.08) (2019.84,4.30)
        (2020.14,3.92) (2021.34,4.20) (2021.54,3.84) (2022,4.10)
        (2022.18,3.78) (2022.36,4.30) (2022.54,3.96)
      };
    \end{axis}
  \end{tikzpicture}

  \vspace{0.8em}

  \begin{tikzpicture}
    \begin{axis}[
      acs xbar,
      width=0.62\linewidth,
      height=4.6cm,
      xmax=16,
      xlabel={Number of publications},
      symbolic y coords={Mechanism/Safety,Material/Surface,Delivery/Combination,Theranostics/Immuno/Other},
      ytick=data,
      y tick label style={font=\scriptsize},
      title={\textbf{(B)} Publications per category}
    ]
      \addplot+[fill=acsblue, draw=none] coordinates {(6,Mechanism/Safety) (8,Material/Surface) (14,Delivery/Combination) (15,Theranostics/Immuno/Other)};
    \end{axis}
  \end{tikzpicture}
  \caption{Publication year point cloud and thematic category distribution of the 43 source publications in the curated selenium nanoparticle dataset.}
  \label{fig:senp-literature-cloud}
\end{figure}

\paragraph{Data Preprocessing.}
For the generalized benchmark dataset, we adhered to the preprocessing protocol established by Masarkar et al. \cite{masarkar2024}:
First, feature columns exhibiting excessive missingness (such as polydispersity index [PDI] and aspect ratio), as well as bibliometric identifiers irrelevant to cytotoxicity prediction (e.g., article ID, publication year, and DOI), were removed, reducing the feature space from an initial 30 descriptors to 15.
Second, nanoparticle classes suffering from severe systemic missingness across fundamental physicochemical parameters (e.g., size, zeta potential, surface charge, and PDI)—specifically quantum dots (QDs) and hydroxyapatite (HAP)—were excluded in their entirety.

Subsequently, the remaining records were harmonized into four overarching nanomaterial classes: metals (M), metal oxides (MO), polymers (P), and carbon-based (C) nanoparticles.
To integrate physicochemical measurements acquired under heterogeneous dispersion conditions, a hierarchical prioritization protocol was implemented:

\begin{enumerate}
  \item \textbf{Size descriptors}:
Measurements obtained in the absence of dispersants (dry/unmodified state) took precedence. If absent, measurements in aqueous dispersion were adopted over cell culture media; if only a single dispersion condition was reported, that solitary value was retained.

  \item \textbf{Zeta potential descriptors}:
Similarly, zeta potential quantified in aqueous suspension was prioritized; if only a single condition was documented, that value was adopted directly.
\end{enumerate}

Following this integration, verification of target endpoint completeness, and systematic anomaly filtering, the final cleaned benchmark dataset retained 4,102 valid data points across 15 feature dimensions.
Continuous descriptors were normalized via min-max scaling ($\text{MinMaxScaler}$), and categorical variables were encoded using label encoding to prepare the tabular matrix for model ingestion.

For the curated SeNP domain-specific dataset, preprocessing procedures were refined to accommodate material-specific nuances.
The resulting SeNP dataset encompasses 383 validated samples and 37 features, spanning realistic operational ranges of exposure concentrations, particle dimensions, morphological types, and surface capping agents, thereby reflecting multi-dimensional relationships between selenium nanoparticles and their physicochemical properties.

The detailed preprocessing protocol for the curated dataset comprised the following stages:

\begin{enumerate}
  \item \textbf{Handling missing values and endpoint filtering}: Because fundamental fields—specifically core particle size, zeta potential, exposure duration, nominal concentration, and cell viability—contained scattered missing entries, records with missing values in these core determinants were eliminated. Furthermore, non-viability entries in the assay method column (e.g., ``Characterization'', ``Western'', and ``Cellular Uptake'') were purged. In total, 262 non-compliant records were eliminated, yielding 383 validated entries.
  \textbf{Note that this row deletion applied strictly to the essential fields specified above; missing values in derived features (e.g., size ratio) and non-critical categorical descriptors (e.g., surface charge) were imputed downstream rather than pruned.}

  \item \textbf{Condition-based priority selection}: In alignment with Masarkar et al. \cite{masarkar2024}, when multiple size measurements were reported for a single formulation, values were selected according to the hierarchy: dry/neat state $>$ aqueous dispersion $>$ cell culture media. Zeta potential was preferentially sourced from aqueous measurements, defaulting to the sole reported condition if alternatives were absent.

  \item \textbf{Dimensional harmonization and unit conversion}: All exposure concentrations were standardized to $\mu$g/mL; concentrations reported in $\mu$M were converted using the atomic mass of selenium (78.96 g/mol). Exposure durations were uniformly transformed into hours (h).

  \item \textbf{Logarithmic feature transformation and ratio derivation}: Given that concentration and exposure time exhibited heavily skewed distributions, both continuous variables underwent log-transformation:
        \[
          \text{Log\_Conc} = \ln(1 + \text{Concentration}_{\mu\text{g/mL}}), \quad
          \text{Log\_Time} = \ln(1 + \text{Time}_{\text{h}})
        \]
        Concurrently, a derived feature—``Size Ratio'' = hydrodynamic diameter / core size—was formulated to capture colloidal agglomeration propensity in aqueous media; missing values in this derived ratio were imputed with a baseline ratio of 1.5.

  \item \textbf{Categorical encoding and imputation}: Categorical attributes including surface charge (positive, negative, neutral), cell lineage, surface functional group, assay type, and crystalline phase (crystalline, amorphous) were label-encoded into integer representations, with missing surface charge categories explicitly imputed as ``Unknown''.

  \item \textbf{Collinearity diagnostics}: Pearson correlation coefficients calculated across all continuous feature pairs exhibited no instances of severe collinearity ($|r| > 0.8$), confirming minimal feature redundancy.
\end{enumerate}

\paragraph{Models and Evaluation.}

On this basis, model input descriptors were stratified into two distinct feature tiers: baseline features and enhanced features. Baseline features encompassed exposure regimes (concentration, duration), foundational physicochemical descriptors (particle core size, zeta potential, size ratio, surface charge), biological host attributes (cell lineage), and experimental parameters (assay method, functional ligand).
For dedicated SeNP modeling, five enhanced descriptors were introduced beyond the baseline set: IC50, elemental Se content, long-term colloidal stability, core-level Se 3d binding energy, and crystalline phase encoding. Among these, elemental Se content, stability, Se 3d binding energy, and crystalline phase directly characterize intrinsic material chemistry. In contrast, because IC50 exhibits an inherent codependency with the cytotoxic endpoint, it was treated as an auxiliary diagnostic descriptor to assess model behavior under endpoint-informed versus purely predictive settings, rather than being conflated with a priori physicochemical descriptors available prior to biological testing.

To systematically quantify the performance gains afforded by the enhanced descriptors while holding baseline features constant, an exhaustive combinatorial search across all $2^5 - 1 = 31$ non-empty subsets was performed. Three-fold cross-validation MAE served as the objective metric to identify optimal feature subsets. Crucially, if an optimal combination incorporated IC50, its predictive gain was interpreted strictly as an auxiliary benchmark rather than an operational performance metric for prospective virtual screening.

We evaluated three paradigm classes for tabular nanotoxicity modeling: standard tree ensembles, attention-augmented tree architectures, and deep tabular neural networks. On the generalized benchmark dataset, six architectures were evaluated: Random Forest, XGBoost, Attention-Based Random Forest (ABRF), Leaf Attention-Based Random Forest (LARF), TabTransformer, and distilled TabTransformer, systematically assessing architectural adaptation to heterogeneous nanoparticle records. On the dedicated SeNP dataset, empirical comparisons focused on performance disparities between baseline and chemically enhanced feature configurations. Random Forest served as the core ensemble baseline to quantify robustness and interpretability under moderate sample regimes; XGBoost was implemented to benchmark the predictive ceiling of gradient-boosted decision trees under enhanced feature configurations; ABRF and LARF were deployed to determine whether hierarchical attention mechanisms improve tree aggregation; and TabTransformer alongside distilled TabTransformer served as deep tabular baselines compared against decision tree paradigms.

Model efficacy was appraised via both regression and classification metrics. In continuous regression tasks, performance was quantified using the coefficient of determination ($R^2$), mean absolute error (MAE), and root mean square error (RMSE). $R^2$ measures the proportion of variance in the target variable explained by the model, bounded from above by 1 (and potentially negative when underperforming an intercept-only mean baseline), with values approaching 1 denoting superior goodness-of-fit. MAE computes the arithmetic average of absolute prediction residuals, sharing the percentage unit (\%) of cellular viability and providing a direct measure of typical prediction discrepancy. RMSE applies quadratic weighting to residuals, exhibiting heightened sensitivity to outlier errors and serving as an indicator of variance stability.

In binary classification tasks, an 85\% cell viability threshold separated samples into ``high viability'' versus ``low viability'' regimes, evaluated using classification accuracy, F1-score, and the area under the receiver operating characteristic curve (AUC). Accuracy denotes the overall proportion of correct classifications; F1-score represents the harmonic mean of precision and recall, exhibiting robustness against class imbalance; and AUC quantifies the discriminative capacity across classification thresholds, where 0.5 denotes random classification and 1.0 reflects ideal separation.

For the generalized multi-type nanoparticle benchmark, a 70:30 train--test split was implemented, with holdout test performance serving as the primary evaluation criterion. Regression tasks were evaluated using $R^2$, MAE, and RMSE, while classification tasks were assessed via Accuracy, F1, and AUC.

For the dedicated SeNP dataset, a three-fold cross-validation scheme was adopted, utilizing CV MAE and overall $R^2$ to assess baseline versus enhanced feature configurations. The automated feature selection phase similarly utilized three-fold cross-validation MAE as the objective function to guarantee cross-subset comparability.

To systematically evaluate the applicability of diverse architectures in nanotoxicity prediction, multiple models were evaluated, spanning conventional tree ensembles, gradient-boosted decision trees, attention-augmented tree models, and deep tabular architectures.
Following initial performance comparisons, six architectures—Random Forest, XGBoost, ABRF, LARF, TabTransformer, and distilled TabTransformer—were selected as the core subjects of analysis.
The architectural design and task-specific positioning of each model are detailed below.

Random Forest is an ensemble learning paradigm built upon decision trees that constructs multiple trees via bootstrap aggregation (bagging) and aggregates individual tree outputs for regression or classification. It exhibits substantial resistance to overfitting, handles mixed feature types (continuous and categorical) seamlessly, shows robustness to missing data, and provides direct feature importance estimations, rendering it an established baseline for small-to-medium tabular tasks.
Accordingly, Random Forest was adopted herein to benchmark model robustness and interpretability on the sample-constrained selenium nanoparticle dataset.

XGBoost is an ensemble framework based on gradient-boosted decision trees that incorporates second-order Taylor expansions of the loss function alongside L1 and L2 regularization penalties, achieving competitive predictive accuracy on tabular benchmarks. By iteratively fitting pseudo-residuals in a sequential manner, XGBoost typically attains high precision, though its sensitivity to hyperparameter configurations necessitates systematic tuning to mitigate overfitting risks.
Herein, XGBoost was introduced to analyze the capacity of gradient boosting when integrating enhanced material-chemistry descriptors (elemental composition, stability, binding energy) and to dissect performance disparities in the presence versus absence of IC50.

ABRF is an ensemble framework that introduces an attention mechanism into standard Random Forest architectures. It first trains a conventional Random Forest and records the predictions generated by individual decision trees across all instances; subsequently, a lightweight neural network dynamically learns instance-specific attention weights for each tree; finally, the ensemble prediction is generated via an attention-weighted summation of tree outputs. The core principle lies in modeling instance-level variations in tree reliability via attention mechanisms to augment ensemble representational capacity.
ABRF was evaluated to assess the impact of single-level attention weighting on tree ensemble predictive fidelity.

LARF builds upon ABRF by introducing an additional feature-level attention mechanism. During model training, it first optimizes a feature attention weight vector to scale input descriptors, thereby prioritizing salient features during tree split selections; it then couples this with tree-level attention to dynamically weight tree predictions, forming a hierarchical, two-level attention architecture. This formulation simultaneously optimizes feature representations and ensemble aggregation.
LARF was implemented to examine whether more complex, multi-level attention mechanisms yield performance gains on the SeNP dataset.

TabTransformer is a deep learning architecture designed specifically for tabular data. Its core principle consists of encoding categorical features into continuous embedding vectors through Transformer self-attention layers, concatenating these contextual embeddings with continuous numerical features, and passing the combined representation through a multi-layer perceptron (MLP) for inference. Because Transformer blocks resolve high-order non-linear interactions among categorical tokens, TabTransformer exhibits pronounced advantages in categorical-heavy tabular settings. Herein, TabTransformer was designated as the deep tabular benchmark to assess deep learning efficacy relative to conventional tree models in nanotoxicity prediction.

Distilled TabTransformer applies knowledge distillation by utilizing the fully trained TabTransformer as a teacher model to train a compact, lightweight student network. By reducing categorical embedding dimensions and narrowing hidden layer widths within the MLP, this framework achieves architectural compression while optimizing against a combined objective of soft teacher targets and ground-truth labels to balance predictive precision and computational efficiency.
This distilled variant was evaluated to appraise the operational feasibility and performance trade-offs of lightweight deep tabular models in resource-constrained deployment environments.

Because all six architectures underwent systematic evaluation on the large-scale, multi-type nanoparticle benchmark, our analysis focuses on architectural adaptation to heterogeneous data distributions.
On the dedicated SeNP dataset, given its constrained sample size (383 records), empirical comparisons centered on three representative architectures: Random Forest, XGBoost, and TabTransformer, assessing performance under baseline versus enhanced feature sets. Random Forest and XGBoost benchmarked standard tree ensembles, while TabTransformer evaluated deep tabular applicability in low-data regimes.
Due to their architectural complexity and elevated computational training overhead, ABRF, LARF, and distilled TabTransformer were evaluated exclusively on the large-scale benchmark to balance experimental complexity, computational resource demands, and analytical interpretability.

\paragraph{Visualization of Feature Effects and Interactions.}
While feature importance metrics identify variables that strongly influence model predictions, ranking alone cannot delineate the functional form of structure--activity relationships—such as whether a response is approximately linear, monotonic, or characterized by thresholded non-linear behavior. Following the analytical framework of Masarkar et al. \cite{masarkar2024}, this study first visualized sample composition, continuous distributions, and categorical frequencies to inspect data coverage and potential selection biases prior to model training. Following model fitting, individual feature effects and pairwise interactions were systematically interrogated using model-agnostic interpretability tools. For univariate effects, partial dependence plots (PDPs) and individual conditional expectation (ICE) curves were utilized to observe average predictive trends alongside instance-level heterogeneity. For bivariate joint effects, two-dimensional PDP contour surfaces were constructed to reveal combinatorial effects on predicted viability or toxicity classes.

Prior to visualizing feature effects, sample composition and variable distributions must be established. The curated SeNP dataset contains 383 validated records (following the exclusion of entries lacking critical physicochemical properties or viability measurements), compiled from 43 peer-reviewed publications spanning 2008--2022. Curated records were concentrated predominantly in 2018 (157 entries, 41.0\%), followed by 2019 (96 entries, 25.1\%) and 2012 (44 entries, 11.5\%). All selenium nanoparticles exhibited spherical morphology, and the biological targets were predominantly derived from human cell lines (371 entries, 96.9\%).

Figure~\ref{fig:senp-year-distribution} illustrates the temporal concentration of source literature, while Table~\ref{tab:senp_feature_statistics} summarizes descriptive statistics for primary continuous features, providing the foundation for subsequent interpretative analyses. Core particle size exhibits pronounced positive skewness (skewness = 3.92), with the majority of particles clustering around 30 nm and rarely exceeding 200 nm, consistent with standard colloidal SeNP synthesis regimes. Hydrodynamic diameter displays milder right-skewness (skewness = 1.69), with its mean (105 nm) substantially exceeding mean core size (65 nm), reflecting the presence of hydration layers and weak colloidal agglomeration. Zeta potential exhibits a near-symmetric distribution (skewness = $-0.08$) ranging from $-60$ mV to $+36$ mV, where negative potentials predominate due to standard surface stabilizers (e.g., polysaccharides, PVP).

\begin{figure}[H]
  \centering
  \begin{minipage}{0.48\linewidth}
    \centering
    \begin{tikzpicture}
      \begin{axis}[
        acs xbar,
        width=\linewidth,
        height=5.1cm,
        xmax=175,
        xlabel={Sample count},
        symbolic y coords={2012,Other years,2019,2018},
        ytick=data,
        y tick label style={font=\scriptsize},
        title={\textbf{(A)} Publication year frequency}
      ]
        \addplot+[fill=acsblue, draw=none] coordinates {(44,2012) (86,Other years) (96,2019) (157,2018)};
      \end{axis}
    \end{tikzpicture}
  \end{minipage}
  \hfill
  \begin{minipage}{0.48\linewidth}
    \centering
    \begin{tikzpicture}
      \begin{axis}[
        acs axis,
        width=\linewidth,
        height=5.1cm,
        ymin=0,
        ymax=105,
        ylabel={Cumulative proportion (\%)},
        symbolic x coords={2018,2019,Other years,2012},
        xtick=data,
        x tick label style={rotate=35, anchor=east, font=\scriptsize},
        ymajorgrids=true,
        title={\textbf{(B)} Cumulative contribution}
      ]
        \addplot+[acsred, mark=square*, mark size=1.7pt] coordinates {(2018,41.0) (2019,66.1) (Other years,88.6) (2012,100.0)};
      \end{axis}
    \end{tikzpicture}
  \end{minipage}
  \caption{Sample size distribution and cumulative contribution across source publication years in the selenium nanoparticle dataset.}
  \label{fig:senp-year-distribution}
\end{figure}

Nominal exposure concentration exhibits extreme right-skewness (skewness = 8.28); while the interquartile range spans $0.52$--$10$ $\mu$g/mL, the maximum concentration reaches $490$ $\mu$g/mL, indicating that a minority of assays employed very high doses that could exert high leverage during model fitting. Exposure duration displays slight negative skewness ($-0.42$) with a median of 72 h, reflecting the prevalence of standard 24 h and 72 h experimental incubation protocols. Cell viability exhibits a near-uniform distribution (skewness = $-0.17$), with both mean and median approaching 50\%, indicating a well-balanced distribution between cytotoxic and non-cytotoxic responses favorable for classification modeling.

To visually synthesize these distributional properties, a coverage-versus-skewness diagnostic plot was constructed for continuous features (Figure~\ref{fig:senp-continuous-summary}). Rather than extrapolating unobserved distributions beyond empirical quantiles, this diagnostic maps available sample counts ($n$) against skewness values from Table~\ref{tab:senp_feature_statistics} to characterize data completeness and distributional distortion.

\begin{table}[H]
  \centering
  \caption{Descriptive statistics of continuous features in the curated SeNP dataset (available sample sizes indicated in the $n$ column)}
  \label{tab:senp_feature_statistics}
  \begin{tabular}{lccccccc}
    \hline
    Feature & $n$ & Mean & SD & Median & Q1 & Q3 & Skewness \\
    \hline
    Core size (nm)               & 339 & 65.30  & 79.72  & 30.00  & 30.00  & 92.15   & 3.92  \\
    Hydrodynamic diameter (nm)   & 279 & 104.97 & 46.73  & 91.00  & 84.00  & 110.00  & 1.69  \\
    Zeta potential (mV)          & 144 & $-5.19$ & 29.76  & $-7.50$ & $-29.10$ & 29.10 & $-0.08$ \\
    Exposure duration (h)        & 363 & 55.21  & 26.14  & 72.00  & 24.00  & 72.00   & $-0.42$ \\
    Concentration ($\mu$g/mL)    & 363 & 12.59  & 38.17  & 0.90   & 0.52   & 10.00   & 8.28  \\
    Cell viability (\%)          & 371 & 49.43  & 28.50  & 50.00  & 41.00  & 68.00   & $-0.17$ \\
    \hline
  \end{tabular}
\end{table}

\begin{figure}[H]
  \centering
  \begin{minipage}{0.49\linewidth}
    \centering
    \begin{tikzpicture}
      \begin{axis}[
        acs ybar,
        width=\linewidth,
        height=5.8cm,
        ymin=0,
        ymax=405,
        ylabel={Available sample size $n$},
        symbolic x coords={Core Size,Hydrodynamic,Zeta,Duration,Concentration,Viability},
        xtick=data,
        x tick label style={rotate=45, anchor=east, font=\scriptsize},
        title={\textbf{(A)} Data coverage}
      ]
        \addplot+[fill=acsblue, draw=none] coordinates {(Core Size,339) (Hydrodynamic,279) (Zeta,144) (Duration,363) (Concentration,363) (Viability,371)};
      \end{axis}
    \end{tikzpicture}
  \end{minipage}
  \hfill
  \begin{minipage}{0.49\linewidth}
    \centering
    \begin{tikzpicture}
      \begin{axis}[
        acs axis,
        width=\linewidth,
        height=5.8cm,
        ymin=-1,
        ymax=9.2,
        ylabel={Skewness},
        symbolic x coords={Core Size,Hydrodynamic,Zeta,Duration,Concentration,Viability},
        xtick=data,
        x tick label style={rotate=45, anchor=east, font=\scriptsize},
        title={\textbf{(B)} Distribution skewness}
      ]
        \addplot+[only marks, acsred, mark=square*, mark size=1.8pt] coordinates {(Core Size,3.92) (Hydrodynamic,1.69) (Zeta,-0.08) (Duration,-0.42) (Concentration,8.28) (Viability,-0.17)};
        \addplot+[black!55, dashed, mark=none] coordinates {(Core Size,0) (Viability,0)};
      \end{axis}
    \end{tikzpicture}
  \end{minipage}

  \vspace{0.9em}

  \begin{minipage}{0.62\linewidth}
    \centering
    \begin{tikzpicture}
      \begin{axis}[
        acs axis,
        width=\linewidth,
        height=5.2cm,
        xmin=110,
        xmax=390,
        ymin=-1,
        ymax=9.2,
        xlabel={Available sample size $n$},
        ylabel={Skewness},
        title={\textbf{(C)} Coverage--skewness diagnostic}
      ]
        \addplot+[
          only marks,
          mark=*,
          mark size=1.9pt,
          acsorange,
          nodes near coords,
          nodes near coords style={font=\tiny, anchor=west, xshift=1.5pt},
          point meta=explicit symbolic
        ] coordinates {
          (339,3.92) [Core Size]
          (279,1.69) [Hydrodynamic]
          (144,-0.08) [Zeta]
          (363,-0.42) [Duration]
          (363,8.28) [Concentration]
          (371,-0.17) [Viability]
        };
        \addplot+[black!55, dashed, mark=none] coordinates {(110,0) (390,0)};
      \end{axis}
    \end{tikzpicture}
  \end{minipage}
  \caption{Available sample size and skewness distribution of primary continuous features.}
  \label{fig:senp-continuous-summary}
\end{figure}

Categorical frequency distributions likewise contextualize the dominant operational domains learned by predictive models (Figure~\ref{fig:senp-categorical-frequency}). Across the investigated dataset, 42 distinct cell lines are documented. Human colorectal carcinoma HCT 116 cells exhibit the highest representation (98 entries, 25.6\%), followed by hepatocellular carcinoma HepG2 (41 entries, 10.7\%), melanoma A375 (23 entries, 6.0\%), and breast adenocarcinoma MCF-7 (23 entries, 6.0\%); the top 10 cell lines collectively encompass 71.0\% of all records. Organ lineage is dominated by colorectal (28.7\%), hepatic (16.7\%), and cerebral (8.9\%) tissues.

Forty-three surface functionalization categories were identified, among which plain polysaccharide coatings (PTR) accounted for the largest share (127 entries, 33.2\%), followed by polysaccharide--protein conjugates (PR; 18 entries, 4.7\%) and PAMAM+RGD+siRNA multifunctional coatings (15 entries, 3.9\%). Synthetic methodologies were heavily centered on chemical reduction routes (352 entries, 91.9\%), while MTT colorimetric reduction represented the predominant cytotoxicity assay format (206 entries, 53.8\%).

\begin{figure}[H]
  \centering
  \begin{minipage}{0.49\linewidth}
    \centering
    \begin{tikzpicture}
      \begin{axis}[
        acs xbar,
        width=\linewidth,
        height=5.8cm,
        xmax=215,
        xlabel={Sample count},
        symbolic y coords={Other cell lines,MCF-7,A375,HepG2,HCT 116},
        ytick=data,
        y tick label style={font=\scriptsize},
        title={\textbf{(A)} Primary cell lines}
      ]
        \addplot+[fill=acsblue, draw=none] coordinates {(198,Other cell lines) (23,MCF-7) (23,A375) (41,HepG2) (98,HCT 116)};
      \end{axis}
    \end{tikzpicture}
  \end{minipage}
  \hfill
  \begin{minipage}{0.49\linewidth}
    \centering
    \begin{tikzpicture}
      \begin{axis}[
        acs axis,
        width=\linewidth,
        height=5.8cm,
        ymin=0,
        ymax=65,
        ylabel={Sample proportion (\%)},
        symbolic x coords={PTR,PR,PAMAM+RGD+siRNA,Other modifications},
        xtick=data,
        x tick label style={rotate=35, anchor=east, font=\tiny},
        nodes near coords,
        nodes near coords style={font=\tiny, yshift=1pt},
        title={\textbf{(B)} Surface functionalization types}
      ]
        \addplot+[only marks, acsorange, mark=square*, mark size=2pt] coordinates {(PTR,33.2) (PR,4.7) (PAMAM+RGD+siRNA,3.9) (Other modifications,58.2)};
      \end{axis}
    \end{tikzpicture}
  \end{minipage}
  \caption{Frequency distribution of primary cell lines and surface functionalization categories.}
  \label{fig:senp-categorical-frequency}
\end{figure}

Consequently, feature visualization within this study was not conducted merely for descriptive profiling, but was integrated into the feature screening and model appraisal workflow: statistical distributions and frequency tables were first leveraged to inspect domain coverage, extreme outlier doses, and class imbalances; post-training interpretability analyses (PDPs, ICE profiles, and two-dimensional contour interactions) were subsequently deployed to interrogate the individual and combinatorial effects of exposure concentration, particle size, zeta potential, incubation duration, assay format, and enhanced chemical descriptors on cytotoxicity predictions.
